\chapter{The Desktop}
\label{ch:desktop}

\standfirst{The window system, in full. It is small, it is consistent, and
none of it works the way anything built after 1990 does.}

\section{What is on the screen}

A grey halftone background, and windows. Each window has a label tab in its
top-left corner and no other decoration: no title bar, no close box, no resize
grip, no scroll bar you have not asked for.

Windows are \emph{scheduled}: \ct{ScheduledControllers} holds the list, and
control passes to whichever window's rectangle contains the cursor. There is
no focus to click into. Move the pointer over a window and that window is
listening.

\section{The three buttons, again}

The colours the library calls them by are in Section~\ref{sec:buttons}; this
is what they do.

\begin{center}
\small
\begin{tabular}{@{}lp{0.62\textwidth}@{}}
\toprule
left & select\\
middle & the view's menu --- what you can do to the contents\\
right & the window's menu --- what you can do to the window\\
\bottomrule
\end{tabular}
\end{center}

\textbf{Press and hold} for both menus. Drag onto an item; release to choose
it; release outside to choose nothing.

\section{The desktop menu}

The middle button on the background:

\begin{center}
\small
\begin{tabular}{@{}lp{0.62\textwidth}@{}}
\toprule
\ctp{restore display} & redraw everything. The cure for a corrupted screen\\
\ctp{exit project} & leave the current project\\
\midrule
\ctp{project} & open a Project: a separate desktop with its own windows\\
\ctp{file list} & a file browser\\
\ctp{browser} & the System Browser\\
\ctp{workspace} & a scratch pad\\
\ctp{system transcript} & where \ctp{Transcript show:} goes\\
\ctp{system workspace} & a workspace pre-filled with useful expressions\\
\midrule
\ctp{save} & write the image, keep running\\
\ctp{quit} & save and quit, quit without saving, or continue\\
\bottomrule
\end{tabular}
\end{center}

Every one of these that opens a window then asks you to frame it: press the
left button,
drag out a rectangle, release.

\section{The window menu}

The right button, anywhere over a window:

\begin{center}
\small
\begin{tabular}{@{}lp{0.62\textwidth}@{}}
\toprule
\ctp{under} & send this window behind the ones it overlaps\\
\ctp{move} & pick it up. It then follows the pointer \emph{with no button
  held}, and you click to drop it\\
\ctp{frame} & give it a new rectangle, the way opening did\\
\ctp{collapse} & shrink it to its label tab, which still has this menu\\
\ctp{close} & close it\\
\bottomrule
\end{tabular}
\end{center}

\ct{close} asks the model first, so a Browser with unaccepted edits will
interrupt and ask whether you meant it.

\section{The text editor}

Every text pane in the system is the same editor, and it has the same
middle-button menu:

\begin{center}
\small
\begin{tabular}{@{}lp{0.60\textwidth}@{}}
\toprule
\ctp{again} & repeat the last replacement, at the next occurrence\\
\ctp{undo} & undo the last change. One level\\
\midrule
\ctp{copy} \ctp{cut} \ctp{paste} & one system-wide clipboard\\
\midrule
\ctp{do it} & evaluate the selection, discard the answer\\
\ctp{print it} & evaluate the selection and insert the answer after it\\
\midrule
\ctp{accept} & commit --- compile the method, save the text\\
\ctp{cancel} & throw the edits away and reload\\
\bottomrule
\end{tabular}
\end{center}

In a Browser's code pane there are three more: \ct{format} (pretty-print),
\ct{spawn} (open a new browser on just this method) and \ct{explain} (say what
the selected token is).

\subsection{Selecting}

Press the left button and drag. Two shortcuts worth having:

\begin{itemize}
\item Double-clicking inside a pair of delimiters---quotes, brackets,
  parentheses---selects everything between them. This is how you select a
  whole block or a whole string quickly.
\item Clicking just inside the leading delimiter and dragging selects to the
  matching one.
\end{itemize}

\subsection{No keyboard shortcuts}

There is no \ct{Ctrl-D}, no \ct{Ctrl-S}, no \ct{Ctrl-Z}. \ct{do it},
\ct{accept} and \ct{undo} are menu items. The only control keys bound are a
few editing and emphasis ones: \ct{Ctrl-W} deletes the previous word,
\ct{Ctrl-B} bolds, \ct{Ctrl-I} italicises, \ct{Ctrl-X} clears emphasis.

\begin{stnote}
Backspace deletes; the arrow keys move the insertion point; \ct{Return}
inserts a carriage return. Typing with a selection active replaces the
selection.
\end{stnote}

\section{Accept means compile}

In a Browser's code pane, \ct{accept} compiles the text as a method of the
selected class and installs it \emph{immediately}. There is no build step,
nothing to restart, and no way to have a stale version running.

If it does not compile you get a \ct{SyntaxError} window with the error marked
in the text. Fix it there and \ct{accept} again, or \ct{cancel}.

\begin{stwarn}
A window whose text you have edited but not accepted is holding your only copy
of that text. Closing it asks first, but \ct{cancel} does not. Accept early.
\end{stwarn}

\section{Projects}

\ct{project} on the desktop menu opens a \ct{Project}: a separate desktop with
its own window list and its own changes. Enter one and the screen is empty
again; \ct{exit project} comes back. It is a virtual desktop, from 1983, and
it is useful when you want a clean screen for one task without closing what
you had.

\section{When the screen goes wrong}

\ct{restore display} on the desktop menu redraws everything from
\ct{ScheduledControllers}. Because the display is an ordinary Form that
anything can draw on---including your own code from a workspace---this is a
normal thing to need, not a sign of a fault.

\section{Scripting the interface}

For tests and screenshots, input can be injected from the command line rather
than typed:

\begin{sh}
./st80 -inject "m 320 240; w 20; d 129; w 40; m 320 271; w 40; u 129; w 60;
                m 60 60; w 20; d 130; w 20; m 900 1400; w 20; u 130; w 200" \
       -screenshot /tmp/browser.pbm -run st80.image
\end{sh}

\ct{m X Y} moves the pointer, \ct{d} and \ct{u} press and release a button
(128 right, 129 middle, 130 left), \ct{k} types a key, \ct{w N} waits N bytecode
slices so the image has time to respond, and \ct{x} posts a window-exposed
event. That script opens the desktop menu, chooses \ct{browser}, frames a
window, and screenshots the result.
